\documentclass[11pt]{article}

% cs250preamble.tex --- shared preamble for CS 250 course notes
%
% This file is \input by main.tex and by each individual module
% file so that each module can still compile standalone. Any
% package or custom-environment change should be made here, not
% in the individual module files.

\usepackage[T1]{fontenc}
\usepackage[margin=1in]{geometry}
\usepackage{amsmath, amssymb, amsthm}
\usepackage{enumitem}
\usepackage{hyperref}
\usepackage{booktabs}
\usepackage{listings}
\usepackage{xcolor}
\usepackage{tikz}
\usetikzlibrary{shapes.geometric, shapes.gates.logic.US, shapes.gates.logic.IEC, arrows.meta, positioning, calc, trees}
\usepackage{bussproofs}
\usepackage{mdframed}

\lstset{
  basicstyle=\ttfamily\small,
  frame=single,
  breaklines=true,
  columns=fullflexible,
  mathescape=true
}

\theoremstyle{definition}
\newtheorem{theorem}{Theorem}[section]
\newtheorem{definition}[theorem]{Definition}
\newtheorem{example}[theorem]{Example}

\hypersetup{
  colorlinks=true,
  linkcolor=blue!60!black,
  urlcolor=blue!60!black
}

% Key-takeaway box: used at the end of major sections and at the end of
% each module to summarise the main points. Expects \item bullets inside.
\newmdenv[
  linewidth=0.8pt,
  linecolor=blue!40!black,
  backgroundcolor=blue!3,
  skipabove=8pt,
  skipbelow=8pt,
  innertopmargin=7pt,
  innerbottommargin=7pt,
  innerleftmargin=9pt,
  innerrightmargin=9pt
]{keytakeawaybox}
\newenvironment{keytakeaway}{%
  \begin{keytakeawaybox}%
  \noindent\textbf{Key takeaways.}%
  \begin{itemize}[leftmargin=1.3em, topsep=2pt, itemsep=1pt, label=\textbullet]%
}{%
  \end{itemize}%
  \end{keytakeawaybox}%
}

% Summary/looking-ahead environment: used once at the end of each module
% to wrap up what the module built and point forward.
\newmdenv[
  linewidth=0.8pt,
  linecolor=gray!60,
  backgroundcolor=gray!5,
  skipabove=8pt,
  skipbelow=8pt,
  innertopmargin=7pt,
  innerbottommargin=7pt,
  innerleftmargin=9pt,
  innerrightmargin=9pt
]{summarybox}

% Sidebar environment: used for tangential asides---historical notes,
% pointers to other areas of math/CS, cross-paradigm remarks---that
% enrich the main narrative without being essential to it. Takes one
% argument: the sidebar's title.
\newmdenv[
  linewidth=0.6pt,
  linecolor=gray!60,
  backgroundcolor=gray!4,
  skipabove=8pt,
  skipbelow=8pt,
  innertopmargin=6pt,
  innerbottommargin=6pt,
  innerleftmargin=9pt,
  innerrightmargin=9pt
]{sidebarbox}
\newenvironment{sidebar}[1]{%
  \begin{sidebarbox}%
  \noindent\textbf{Sidebar: #1.}\par\medskip%
}{%
  \end{sidebarbox}%
}

% Reading-map environment: used at the top of each numbered module to
% indicate Epp coverage, prerequisites, relevant intermezzi, and
% suggested practice problems. Expects four \rmentry{label}{body}
% items inside (or arbitrary paragraphs).
\newmdenv[
  linewidth=0.8pt,
  linecolor=black!55,
  backgroundcolor=yellow!4,
  skipabove=10pt,
  skipbelow=14pt,
  innertopmargin=9pt,
  innerbottommargin=9pt,
  innerleftmargin=11pt,
  innerrightmargin=11pt
]{readingmapbox}
\newcommand{\rmentry}[2]{\par\noindent\textbf{#1}\hspace{0.4em}#2\par}
\newenvironment{readingmap}{%
  \begin{readingmapbox}%
  \noindent{\bfseries\large Reading map}%
  \vspace{2pt}%
  \setlength{\parskip}{4pt plus 1pt}%
}{%
  \end{readingmapbox}%
}


\title{CS 250 --- How to Write a Proof}
\author{}
\date{}

\begin{document}
\maketitle

You now have enough vocabulary---sets, functions, relations---to start writing proofs about them. But when you sit down to actually write one, it can be hard to know where to begin. This short chapter gives practical advice about formatting and structure. We are not learning proof \emph{techniques} yet (that starts in Module~3); we're just learning how to put words on the page.

Keep this chapter bookmarked. The specific proof techniques change as the course goes on---direct, counterexample, cases, induction, contradiction, contrapositive---but the prose habits below carry across all of them.

\section{A proof is a piece of writing}

A proof is not a calculation, a formula, or a wall of symbols. It is a short essay that explains, step by step, \emph{why} something is true. Each sentence should either state a fact, introduce a definition, or draw a conclusion from what came before. If you can't explain a step in words, you probably don't understand it yet.

Here are the ground rules:
\begin{enumerate}
  \item \textbf{State what you are proving.} Open with something like ``We prove that \ldots'' or ``We want to show that \ldots''. Don't make the reader guess your goal.
  \item \textbf{Introduce variables before you use them.} If your proof involves a number, a set, or a function, give it a name and say where it comes from. ``Let $n$ be a positive integer'' or ``Consider an arbitrary $x \in \mathbb{R}$'' or ``Define $f : A \to B$ by $f(x) = x^2$.''
  \item \textbf{Write in complete sentences.} Mathematical symbols are parts of sentences, not replacements for them. Don't write a bare chain of equations with no words connecting them.
  \item \textbf{Justify each step.} After every claim, the reader should be able to see \emph{why} it follows---from a definition, from a previous step, or from a well-known fact.
  \item \textbf{End with a clear conclusion.} Finish with ``Therefore \ldots'' or ``This completes the proof.'' The reader should know when you're done and what you've established.
\end{enumerate}

\section{When the problem has cases}

Many proofs, including ones you'll see on this week's assignment, naturally split into cases. For instance, you might need to show something holds ``for all positive integers $A$,'' but the argument is different depending on whether $A$ is even or odd. In that situation:

\begin{itemize}
  \item State explicitly that you are splitting into cases and say what they are: ``We consider two cases: $A$ is even, and $A$ is odd.''
  \item Handle each case separately. Label them clearly (``\textbf{Case 1:} $A$ is even.'') and give a self-contained argument for each one.
  \item After the cases, tie them together: ``Since every positive integer is either even or odd, these two cases are exhaustive, and the proof is complete.''
\end{itemize}

The key requirement is that your cases \emph{cover all the possibilities}. If there's a scenario you didn't address, your proof has a gap.

\section{Using definitions}

When a proof asks whether something is or isn't a function, or whether a relation has a particular property, \emph{go back to the definition}. This is the single most important habit in proof-writing.

For example, if you need to show that a certain relation is a function, you should explicitly invoke the definition: a function $f : A \to B$ assigns every element of $A$ to \emph{exactly one} element of $B$. So you need to show two things: (1) for every input there is at least one output, and (2) for every input there is at most one output. Conversely, to show something is \emph{not} a function, it suffices to find a single input that has zero outputs or more than one.

Don't skip this step. Even if the answer feels obvious, the proof isn't complete until you've connected your argument to the actual definitions.

\section{An example: good vs.\ bad}

Suppose we want to prove that the relation $y^2 = x$ on the reals is not a function (with $x$ as input and $y$ as output).

\medskip
\noindent\textbf{Bad version:}
\begin{quote}
$y^2 = x$ is not a function because $x = 4$ gives $y = 2$ and $y = -2$.
\end{quote}

This has the right idea but it doesn't explain \emph{why} having two values of $y$ is a problem. The reader is left to fill in the reasoning.

\medskip
\noindent\textbf{Good version:}
\begin{quote}
We show that the relation $y^2 = x$ on $\mathbb{R}$ is not a function from $x$ to $y$. By definition, a function must assign each input $x$ to exactly one output $y$. Consider $x = 4$. Then $y^2 = 4$, which has two real solutions: $y = 2$ and $y = -2$. Since the input $x = 4$ corresponds to two different outputs, the relation violates the uniqueness requirement and is therefore not a function.
\end{quote}

Notice what the good version does: it states the goal, invokes the definition, gives the specific example, and explicitly connects the example back to the definition to reach the conclusion. Every step is justified. That's all a proof needs to be.

\section{Where to go from here}

The advice above is independent of any particular proof technique. In the modules that follow you will meet the standard techniques---direct proof (Modules~3 and~5), counterexample and proof by cases (Modules~5 and~6), induction (Modules~4 and~8--10), proof by contradiction and contrapositive (Module~7)---and each will come with its own recipe. But every one of them expects to be written in prose that follows the rules above: state the goal, introduce variables, justify each step, conclude clearly.

Plan to re-read this chapter after every few proofs you write, especially early on. The mistakes most worth fixing are not mathematical---they are writing mistakes: missing variable introductions, unjustified jumps, equations strung together without connective tissue, no clear conclusion. Fix the writing and the mathematics usually becomes visible.

\end{document}
